\chapter{Introducción}
\label{cap:introduccion}

\section{Contexto}
\label{sec:introduccion:contexto}

Este trabajo se ha realizado en el contexto de las líneas de investigación
del grupo ISA (Ingeniería del Software Aplicada) de la Universidad de
Sevilla, y en particular de la línea de \textit{Pricing Intelligence}
\citep{cavero2026pricingintelligence}.

\subsection{Los dos canales de un servicio SaaS}
\label{subsec:introduccion:canales}

Un servicio ofrecido bajo el modelo \textit{Software as a Service} (SaaS)
puede ser consumido a través de distintos canales de interacción. El más
visible es habitualmente un \textit{frontend} web, orientado a usuarios
humanos, en el que las condiciones de contratación se publican de forma
tabular mediante planes, características (\textit{features}) y complementos
(\textit{add-ons}). Junto a él, el mismo servicio puede exponer una API
(REST o GraphQL) para que otros sistemas consuman sus capacidades de
forma programática.

Aunque ambos canales describen aspectos del mismo servicio, cada uno
utiliza fuentes, estructuras y terminologías distintas. El canal web habla
de planes y funcionalidades; el canal API habla de \textit{endpoints},
límites de velocidad (\textit{rates}), cuotas de volumen y reglas de
consumo.
La Sección~\ref{sec:estado:acm} precisa estos términos: con la definición del NIST, lo que aquí se llama canal es una interfaz del servicio.

En el dominio del correo electrónico esta dualidad es especialmente
marcada. Proveedores como SendGrid, Mailgun, Postmark, Resend o Amazon SES
publican una página de precios orientada a la contratación y, en paralelo,
una documentación técnica donde residen las condiciones que realmente
gobiernan el consumo programático. La información determinante para un
integrador no siempre aparece en la página de precios: cuántas peticiones
por segundo admite un endpoint, cuántos destinatarios caben en una
llamada, qué ocurre al superar la cuota. Está
dispersa entre documentación de desarrollador, artículos de soporte,
especificaciones OpenAPI y políticas de uso aceptable. El
Capítulo~\ref{cap:hallazgos} mide, caso por caso, cuánto de esa
información falta.

\subsection{Pricing Intelligence y su ecosistema}
\label{subsec:introduccion:pi}

El \textit{Pricing Intelligence} se define como una disciplina dirigida por
datos cuyo objetivo es transformar la información contenida en los planes
de precios en conocimiento útil para la toma de decisiones
\citep{cavero2026pricingintelligence}. Recopilar precios es solo el
punto de partida. El objetivo es responder de forma automática a preguntas
como qué plan se ajusta mejor a un patrón de consumo dado, o qué
diferencias relevantes existen entre dos proveedores.

Alrededor de esta idea, el grupo ha construido un ecosistema de
herramientas interoperables. Para el canal web, el metamodelo
\textit{iPricing} y su serialización Pricing2YAML representan los planes
de un SaaS \citep{garciafernandez2025ipricing}. Sobre ellos operan A-MINT,
extractor de pricings desde páginas web mediante modelos de lenguaje
\citep{cavero2025amint}, y PRIME, validador y analizador de los modelos
resultantes \citep{garciafernandez2025ipricing}. HARVEY, presentado en
CAiSE 2026, es el orquestador de estas capacidades tras una interfaz
conversacional \citep{cavero2026pricingintelligence}. Finalmente, SPHERE es el
\textit{hub} del ecosistema: ahí se publican los artefactos, ahí se ven y
desde ahí se comparten \citep{sphere}.

La extensión de este ecosistema al canal API es reciente. El Trabajo Fin de
Máster de Gavira Sánchez \citep{gavira2026tfm}, en coordinación con el de
Ruiz López \citep{ruizlopez2026tfm}, conceptualiza la capacidad contractual
de una API en el \textit{API Capacity Model} (ACM) y define ACM2YAML, la
serialización que da lugar a la \textbf{Datasheet}. La Datasheet es un
documento orientado a máquinas que reúne en una sola representación la
información tabular del pricing y la que aparece dispersa en la
documentación técnica. Sobre ella operan PRIME4API \citep{prime4apirepo},
motor de análisis de la capacidad, y HARVEY4API \citep{ruizlopez2026tfm},
el agente que lo consume.

\section{Motivación}
\label{sec:introduccion:motivacion}

El responsable técnico que debe elegir entre una docena de proveedores
rara vez encuentra en la página de precios los tres datos de consumo que
enumera la Sección~\ref{subsec:introduccion:canales}. Se enfrenta, además, a
documentaciones heterogéneas y a comparativas elaboradas por los propios
vendedores \citep{mailtrap2026comparativa,mailslurp2026comparativa}. Y le falta de forma sistemática la información que más le
importa: la que determina si el servicio soportará su patrón de consumo real.

El ecosistema descrito en la sección anterior ofrece el
instrumental necesario para abordar el problema de forma rigurosa.
Sin embargo, ese instrumental se ha validado hasta ahora sobre tres
pricings de dos servicios \citep{gavira2026tfm}, y su versión operativa captura
dos de las cuatro dimensiones de capacidad del metamodelo que la
sustenta. Existe, por tanto, una distancia entre lo que el modelo promete
y lo que se ha demostrado que soporta.

Esa distancia es la que motiva este trabajo, y es también la que determina
su dificultad. Los motores de análisis ya existen, así que el reto no está
en la complejidad algorítmica. Está en tres frentes de naturaleza distinta.
El primero es de reconstrucción. La capacidad contractual de una API real
rara vez está publicada en un único lugar, y en ocasiones sencillamente no
está publicada. Documentarla exige localizar, integrar y contrastar fuentes
heterogéneas y con frecuencia incompletas. El segundo es de modelado. Cada
servicio que no encaja en el esquema obliga a discernir si es un caso
particular que puede simplificarse o una carencia real del modelo que debe
corregirse, y esa distinción no admite automatización. El tercero es de
coordinación. Este trabajo se construye sobre artefactos que otros
trabajos desarrollan en paralelo, y cualquier cambio propuesto debe
verificarse contra el resto del ecosistema.

\section{Objetivos y preguntas de investigación}
\label{sec:introduccion:objetivos}

Del contexto y la motivación expuestos se deriva el objetivo global de este
Trabajo Fin de Máster.

\begin{mdframed}[style=US_style]
\textbf{Objetivo del Trabajo Fin de Máster}\\[.2em]
Construir un espacio de conocimiento sobre las APIs de correo electrónico
mediante Datasheets conformes al ACM, y establecer un método de comparación
sistemática que responda a necesidades funcionales y no funcionales,
evolucionando el modelo con las carencias que el dominio revele.
\end{mdframed}

Este objetivo se articula en torno a cuatro preguntas de investigación.

\begin{description}
  \item[RQ1.] ¿Qué APIs de correo deben incluirse en un análisis
    representativo del dominio, y bajo qué criterios reproducibles?
  \item[RQ2.] ¿Es ACM2YAML suficientemente expresivo para representar la
    capacidad contractual de las APIs de correo? ¿Qué extensiones requiere?
  \item[RQ3.] ¿Cómo se relacionan los límites de uso y los \textit{rates}
    en las APIs de correo, y qué implicaciones tiene esa relación para el
    consumidor?
  \item[RQ4.] ¿Cómo pueden compararse las APIs de correo de forma
    sistemática para responder a preguntas de decisión que involucran
    simultáneamente requisitos funcionales y no funcionales?
\end{description}

De ellas se derivan los siguientes objetivos específicos.

\begin{description}
  \item[OBJ-01] \textbf{Selección justificada del dominio.} Definir y
    aplicar un criterio de inclusión reproducible que determine el conjunto
    de APIs analizadas, documentando el proceso de forma que un tercero
    pueda replicarlo.
  \item[OBJ-02] \textbf{Construcción del catálogo de Datasheets.}
    Documentar la capacidad contractual de cada API seleccionada conforme
    a ACM2YAML, con trazabilidad explícita de las fuentes de cada dato.
  \item[OBJ-03] \textbf{Evolución del esquema.} Identificar las carencias
    de expresividad que el dominio revele y proponer las extensiones
    correspondientes, justificando cada una con el caso real que la motiva.
  \item[OBJ-04] \textbf{Análisis de la capacidad.} Caracterizar cómo se
    relacionan los límites de uso y los \textit{rates} en el dominio, con la
    lectura geométrica de la capacidad y las operaciones de análisis de
    PRIME4API.
  \item[OBJ-05] \textbf{Método de comparación.} Establecer un catálogo de
    preguntas de decisión propias del dominio y un procedimiento para
    responderlas sobre el conjunto de Datasheets.
\end{description}

El componente de investigación del trabajo reside principalmente en OBJ-03
y OBJ-04. La construcción del catálogo es condición necesaria, pero es un
medio. El resultado investigable es lo que la confrontación entre un modelo
y un dominio real revela sobre las limitaciones del primero, y el
conocimiento que puede extraerse del segundo una vez formalizado.

\section{Contribuciones}
\label{sec:introduccion:contribuciones}

El desarrollo de estos objetivos se materializa en cuatro contribuciones.

\begin{description}
  \item[C-01] \textbf{Catálogo de Datasheets de APIs de correo.} Un
    conjunto curado y versionado de Datasheets del dominio, con trazabilidad
    de fuentes y preparado como contribución a SPHERE, reutilizable por el resto de
    herramientas del ecosistema.
  \item[C-02] \textbf{Evolución de ACM2YAML.} Una propuesta de extensión del
    esquema justificada caso a caso, en la que cada elemento nuevo responde
    a un escenario real que la versión previa no era capaz de representar.
  \item[C-03] \textbf{Caracterización de la relación entre límites y
    \textit{rates}.} Un análisis de cómo interactúan ambas dimensiones en el
    dominio del correo y de sus consecuencias para el consumidor. La
    relación estaba definida en el modelo, pero no caracterizada sobre un
    dominio real (Sección~\ref{sec:estado:analisis}).
  \item[C-04] \textbf{Método de comparación sistemática.} Un catálogo de
    preguntas de decisión funcionales y no funcionales, junto con el
    procedimiento para resolverlas sobre las Datasheets.
\end{description}

Estas contribuciones tienen aplicación académica e industrial. En el
ámbito aca\-dé\-mi\-co, el catálogo de Datasheets es un banco de pruebas sobre
un dominio real para las herramientas del ecosistema, que hasta ahora
se han validado sobre tres pricings de dos servicios \citep{gavira2026tfm}.
Las extensiones propuestas al esquema revierten directamente sobre el
modelo. En el ámbito industrial, el instrumento aborda una decisión que
los equipos de desarrollo toman de forma recurrente y sin base
comparable. Cabe esperar que el método y el esquema se extiendan a otros
dominios de APIs más allá del correo electrónico, aunque esa
extensibilidad queda sin probar (Sección~\ref{sec:conclusiones:externa}).

\section{Estructura de la memoria}
\label{sec:introduccion:estructura}

El \textbf{Capítulo~\ref{cap:introduccion}} introduce el contexto del
trabajo, presenta los dos canales de acceso a un SaaS y el ecosistema de
Pricing Intelligence, formula la motivación y los objetivos, y enumera las
contribuciones realizadas.

El \textbf{Capítulo 2} describe la metodología seguida y la planificación
del proyecto, con su estructura de paquetes de trabajo, el presupuesto
estimado y las desviaciones observadas.

El \textbf{Capítulo 3} revisa el estado del arte en el modelado de pricings
de SaaS, en el modelado y análisis de la capacidad de APIs, y en las
herramientas del ecosistema de Pricing Intelligence sobre las que se
construye este trabajo.

El \textbf{Capítulo 4} presenta el método con el que se construye el
catálogo (criterio de selección del dominio, ciclo de construcción por
API, reglas de trazabilidad y confianza) y su calibración sobre un piloto
real.
Cierra la RQ1 con la condición que debe cumplir el conjunto de APIs
del catálogo (Sección~\ref{sec:propuesta:respuesta}).

El \textbf{Capítulo 5} recoge lo que las APIs del catálogo revelan. Reúne
los hallazgos de cada una, la relación observada entre límites de uso y
\textit{rates}, y el catálogo de preguntas de Pricing Intelligence con
el que se comparan de forma sistemática. Las Datasheets completas no se
reproducen en la memoria. Se conservan, con sus fuentes y evidencias, en
el repositorio de datos que acompaña al trabajo.

El \textbf{Capítulo 6} formula la evolución del esquema ACM2YAML. Recoge
el inventario de elementos no soportados levantado sobre las ocho APIs del
catálogo, aplica
los criterios de promoción y propone las extensiones que el dominio
justifica, cada una con el caso real que la motiva.

El \textbf{Capítulo 7} cierra la memoria. Responde a las cuatro preguntas de
investigación y evalúa las contribuciones frente a los objetivos. Después
agrupa las amenazas a la validez y describe el paquete de replicación.
Termina con las líneas de trabajo futuro y con las conclusiones técnicas.

El Anexo~\ref{anx:uso-ia} declara el uso de inteligencia artificial en el trabajo, con la consulta que completó la selección de APIs.
